%% Time-stamp: <2018-10-18 20:24:12 (marc)>
\documentclass[xcolor=x11names,compress,mathserif]{beamer}

\newcommand{\hackspace}{\hspace{4.2mm}}
\newcommand{\showstudent}[1]{}
\newcommand\hmmax{0}
\newcommand\bmmax{0}


\usepackage{../includes/MarkMathCmds}





% talk/author information
\newcommand{\authorname}{Mark van der Wilk}
\newcommand{\authoremail}{m.vdwilk@imperial.ac.uk}
\newcommand{\authoraffiliation}{
  Department of Computing\\Imperial
  College London}
\newcommand{\authortwitter}{markvanderwilk}
\newcommand{\slidesettitle}{\imperialBlue{Vector Calculus}}
\newcommand{\footertitle}{Differentiation}
\newcommand{\location}{Imperial College London}
\newcommand{\talkDate}{October 12, 2021}



\date{\imperialGray{\talkDate}}




% load defaults
\selectcolormodel{rgb}
\usepackage{ifxetex,ifluatex}
\newif\ifxetexorluatex
\ifxetex
  \xetexorluatextrue
\else
  \ifluatex
    \xetexorluatextrue
  \else
    \xetexorluatexfalse
  \fi
\fi

\usepackage{textpos}
%\usepackage{arabtex}
\usepackage{tikz}
\usetikzlibrary{decorations.markings}
\usetikzlibrary{arrows}
\usetikzlibrary{shapes}
\usetikzlibrary{plotmarks}
\usetikzlibrary{mindmap,trees,backgrounds}

\tikzstyle{every picture}+=[remember picture]

%\usepackage{movie15}
% \usepackage{pdfpages}
%\usepackage{xmpmulti}

\usepackage{anyfontsize}
\usepackage{wrapfig}
\usepackage{animate}
\usepackage{multirow}
\usepackage{multimedia}
\usepackage{xmpmulti}
%\usepackage[latin9]{inputenc}
\usepackage[english]{babel}
\usepackage{scalefnt}
\usepackage{verbatim}
\usepackage{url}
% \usepackage{pgf,pgfarrows,pgfnodes}
\usepackage{textpos}
\usepackage[tight,ugly]{units}
\usepackage{url}
\usepackage{bbm}
\usepackage[english]{babel}
\usepackage{fancyhdr}
\usepackage{bm} % correct bold symbols, like \bm
\usepackage{amsmath}
\usepackage{amsfonts}
\usepackage{amssymb}
\usepackage{mathrsfs}
\usepackage{mathtools}
\usepackage{color}
\usepackage{cancel}
\usepackage{algorithm}
\usepackage{algpseudocode}
\usepackage{mathrsfs}
\usepackage{listings}
\usepackage{graphicx} % for pdf, bitmapped graphics files
\usepackage{mathtools}
\usepackage{units}
\usepackage{subfig}
\usepackage{enumerate}
\usepackage[comma,authoryear]{natbib}
\usepackage{dsfont}


\ifxetexorluatex
\usepackage{fontspec}
\setmainfont[Scale=0.8]{OpenDyslexic-Regular}
\else
\usefonttheme{professionalfonts}
\fi

\renewcommand{\vec}[1]{{\boldsymbol{{#1}}}} % vector
\newcommand{\mat}[1]{{\boldsymbol{{#1}}}} % matrix
% \newcommand{\KL}[2]{\mathrm{KL}(#1\|#2)} % KL divergence
\newcommand{\R}[0]{\mathds{R}} % real numbers
\newcommand{\Z}[0]{\mathds{Z}} % integers
\newcommand{\tr}[0]{\text{tr}} % trace
% \newcommand{\inv}{^{-1}}
% \DeclareMathOperator*{\diag}{diag}
\newcommand{\E}{\mathds{E}} % expectation
\newcommand{\var}{\mathds{V}}
\newcommand{\gauss}[2]{\mathcal{N}\big(#1,\,#2\big)}
\newcommand{\gaussx}[3]{\mathcal{N}\big(#1\,|\,#2,\,#3\big)}
\newcommand{\gaussBig}[2]{\mathcal{N}\left(#1,\,#2\right)}
\newcommand{\gaussxBig}[3]{\mathcal{N}\left(#1\,\left|\,#2,\,#3\right.\right)}
\newcommand{\Ber}[0]{\mathrm{Ber}} % Bernoulli distribution
\DeclareMathOperator{\cov}{Cov}
\ifxetexorluatex
\renewcommand{\T}[0]{^\top}
\renewcommand{\d}[0]{\text{d}} % derivative
\else
\newcommand{\T}[0]{^\top}
\renewcommand{\d}[0]{\text{d}} % derivative
\fi
% calculus
\newcommand{\pdiff}[1]{\frac{\partial}{\partial #1}}
\newcommand{\pdiffF}[2]{\frac{\partial #1}{\partial #2}}
\newcommand{\diffF}[2]{\frac{{\d}#1}{{\d}#2}}
\newcommand{\diffFII}[2]{\frac{{\d}^2 #1}{{\d}#2^2}}
\newcommand{\diff}[1]{\frac{{\d}}{{\d}#1}}
\newcommand{\diffII}[1]{\frac{{\d}^2}{{\d}#1^2}}
\newcommand{\class}[0]{\mathcal{C}}

\newcommand{\idx}[1]{{(#1)}}
% \newcommand{\norm}[1]{\left\|#1\right\|}
\newcommand{\proj}[1]{\tilde{#1}}
\newcommand{\pcacoord}{z}
\newcommand{\pcacoordnew}{\zeta}
\newcommand{\latent}{z}
% \newcommand{\given}{\,|\,}
\newcommand{\genset}[1]{\mathrm{span}[#1]} % generating set
\newcommand{\set}[1]{\mathcal{#1}} % set
\newcommand{\fixgmfont}[1]{\scalebox{0.8}{#1}}



\usepackage{pifont}% http://ctan.org/pkg/pifont
\newcommand{\cmark}{{\color{green!40!black}\ding{51}}}%
\newcommand{\xmark}{{\color{red}\ding{55}}}%
\newcommand{\green}[1]{{\bf{\textcolor{green}{#1}}}}
\newcommand{\red}[1]{{\bf{\textcolor{red}{#1}}}}

\newcommand<>\red[1]{{\color#2[rgb]{1,0,0}#1}}
\newcommand<>\blue[1]{{\color#2[rgb]{0,0,1}#1}}
\newcommand<>\yellow[1]{{\color#2{camyellow}#1}}
\newcommand<>\green[1]{{\color#2[rgb]{0,0.6,0.0}#1}}
\newcommand<>\violet[1]{{\color#2[rgb]{0.6,0,0.6}#1}}
\newcommand<>\orange[1]{{\color#2[rgb]{1,0.5,0}#1}}
\newcommand<>\black[1]{{\color#2[rgb]{0,0,0}#1}}
\newcommand<>\steel[1]{{\color#2[rgb]{0,0,0.8}#1}}
\newcommand<>\darkblue[1]{{\color#2[rgb]{0,0,0.6}#1}}
\newcommand<>\lightblue[1]{{\color#2[rgb]{0.4,0.4,0.7}#1}}
\newcommand<>\gray[1]{{\color#2[rgb]{0.4,0.4,0.4}#1}}
\newcommand<>\greenish[1]{{\color#2[rgb]{0.45, 0.66, 0.45}#1}}
\newcommand<>\redish[1]{{\color#2[rgb]{0.7843    0.3706    0.3706}#1}}
\definecolor{redishTIKZ}{rgb}{0.7843, 0.3706, 0.3706}
\definecolor{imperialBlue}{rgb}{0.058, 0.219, 0.418}
\definecolor{aimsbrown}{rgb}{0.539, 0.117, 0.015}
% \definecolor{imperialGray}{rgb}{0.414, 0.488, 0.671 }
\definecolor{imperialGray}{RGB}{109,153, 204}
\definecolor{aimslightbrown}{RGB}{138,88,84}
\newcommand<>\imperialBlue[1]{{\color#2[rgb]{0.058, 0.219, 0.418}#1}}
\newcommand<>\aimsbrown[1]{{\color#2[rgb]{0.539, 0.117, 0.015}#1}}
%\newcommand<>\imperialGray[1]{{\color#2[rgb]{0.414, 0.488, 0.671}#1}}
\newcommand<>\imperialGray[1]{{\color#2[RGB]{109,153, 204}#1}}
\newcommand<>\aimslightbrown[1]{{\color#2[RGB]{138,88,84}#1}}
\newcommand<>\lightgray[1]{{\color#2[rgb]{0.8,0.8,0.8}#1}}
%\newcommand<>\highlightcolor[1]{{\color#2[rgb]{0,0,1}#1}}
\newcommand{\highlight}[1]{{\bf\steel{#1}}}
%\newcommand{\newblock}[0]{}

%\newcommand{\arrow}[0]{\includegraphics[height=5pt]{./figures/arrow}\hspace{3pt}}

\renewcommand{\emph}[1]{\textbf{\steel{{#1}}}}

\renewcommand{\alert}[1]{{\bf\red{{#1}}}}

\newcommand{\arrow}{
\begin{tikzpicture}
\draw [black!40!green, fill=black!40!green] (0,-0.12) -- (0,0.12) --
(0.15,0);
\draw [black!40!green, fill=black!40!green] (0.15,-0.12) -- (0.15,0.12) --
(0.3,0); 
\end{tikzpicture}
}

\geometry{left=0.45cm,top=0cm,right=0.45cm}


\newcommand{\logoimagepath}{./figures/imperial}
\newcommand{\highlightcolor}{blue!80!black}
%\newcommand{\headbarcolor}{imperialBlue}
\newcommand{\headbarcolor}{imperialBlue}
\institute{}

\newcommand{\coursetitle}{}

\newcommand{\slidesetsubtitle}{}
\newcommand{\slidesetnumber}{01}
\usefonttheme{professionalfonts}


\usetikzlibrary{decorations.fractals}
% tikzlibrary.code.tex
%
% Copyright 2010-2011 by Laura Dietz
% Copyright 2012 by Jaakko Luttinen
%
% The MIT License
%
% See LICENSE file for more details.

% Load other libraries
\usetikzlibrary{shapes}
\usetikzlibrary{fit}
\usetikzlibrary{chains}
\usetikzlibrary{arrows}

% Latent node
\tikzstyle{latent} = [circle,fill=white,draw=black,inner sep=1pt,
minimum size=20pt, font=\fontsize{10}{10}\selectfont, node distance=1]
% Observed node
\tikzstyle{obs} = [latent,fill=gray!25]
% Constant node
\tikzstyle{const} = [rectangle, inner sep=0pt, node distance=1]
% Factor node
\tikzstyle{factor} = [rectangle, fill=black,minimum size=5pt, inner
sep=0pt, node distance=0.4]
% Deterministic node
\tikzstyle{det} = [latent, diamond]

% Plate node
\tikzstyle{plate} = [draw, rectangle, rounded corners, fit=#1]
% Invisible wrapper node
\tikzstyle{wrap} = [inner sep=0pt, fit=#1]
% Gate
\tikzstyle{gate} = [draw, rectangle, dashed, fit=#1]

% Caption node
\tikzstyle{caption} = [font=\footnotesize, node distance=0] %
\tikzstyle{plate caption} = [caption, node distance=0, inner sep=0pt,
below left=5pt and 0pt of #1.south east] %
\tikzstyle{factor caption} = [caption] %
\tikzstyle{every label} += [caption] %

%\pgfdeclarelayer{b}
%\pgfdeclarelayer{f}
%\pgfsetlayers{b,main,f}

% \factoredge [options] {inputs} {factors} {outputs}
\newcommand{\factoredge}[4][]{ %
  % Connect all nodes #2 to all nodes #4 via all factors #3.
  \foreach \f in {#3} { %
    \foreach \x in {#2} { %
      \path (\x) edge[-,#1] (\f) ; %
      %\draw[-,#1] (\x) edge[-] (\f) ; %
    } ;
    \foreach \y in {#4} { %
      \path (\f) edge[->, >={triangle 45}, #1] (\y) ; %
      %\draw[->,#1] (\f) -- (\y) ; %
    } ;
  } ;
}

% \edge [options] {inputs} {outputs}
\newcommand{\edge}[3][]{ %
  % Connect all nodes #2 to all nodes #3.
  \foreach \x in {#2} { %
    \foreach \y in {#3} { %
      \path (\x) edge [->, >={triangle 45}, #1] (\y) ;%
      %\draw[->,#1] (\x) -- (\y) ;%
    } ;
  } ;
}

% \factor [options] {name} {caption} {inputs} {outputs}
\newcommand{\factor}[5][]{ %
  % Draw the factor node. Use alias to allow empty names.
  \node[factor, label={[name=#2-caption]#3}, name=#2, #1,
  alias=#2-alias] {} ; %
  % Connect all inputs to outputs via this factor
  \factoredge {#4} {#2-alias} {#5} ; %
}

% \plate [options] {name} {fitlist} {caption}
\newcommand{\plate}[4][]{ %
  \node[wrap=#3] (#2-wrap) {}; %
  \node[plate caption=#2-wrap] (#2-caption) {#4}; %
  \node[plate=(#2-wrap)(#2-caption), #1] (#2) {}; %
}

% \gate [options] {name} {fitlist} {inputs}
\newcommand{\gate}[4][]{ %
  \node[gate=#3, name=#2, #1, alias=#2-alias] {}; %
  \foreach \x in {#4} { %
    \draw [-*,thick] (\x) -- (#2-alias); %
  } ;%
}

% \vgate {name} {fitlist-left} {caption-left} {fitlist-right}
% {caption-right} {inputs}
\newcommand{\vgate}[6]{ %
  % Wrap the left and right parts
  \node[wrap=#2] (#1-left) {}; %
  \node[wrap=#4] (#1-right) {}; %
  % Draw the gate
  \node[gate=(#1-left)(#1-right)] (#1) {}; %
  % Add captions
  \node[caption, below left=of #1.north ] (#1-left-caption)
  {#3}; %
  \node[caption, below right=of #1.north ] (#1-right-caption)
  {#5}; %
  % Draw middle separation
  \draw [-, dashed] (#1.north) -- (#1.south); %
  % Draw inputs
  \foreach \x in {#6} { %
    \draw [-*,thick] (\x) -- (#1); %
  } ;%
}

% \hgate {name} {fitlist-top} {caption-top} {fitlist-bottom}
% {caption-bottom} {inputs}
\newcommand{\hgate}[6]{ %
  % Wrap the left and right parts
  \node[wrap=#2] (#1-top) {}; %
  \node[wrap=#4] (#1-bottom) {}; %
  % Draw the gate
  \node[gate=(#1-top)(#1-bottom)] (#1) {}; %
  % Add captions
  \node[caption, above right=of #1.west ] (#1-top-caption)
  {#3}; %
  \node[caption, below right=of #1.west ] (#1-bottom-caption)
  {#5}; %
  % Draw middle separation
  \draw [-, dashed] (#1.west) -- (#1.east); %
  % Draw inputs
  \foreach \x in {#6} { %
    \draw [-*,thick] (\x) -- (#1); %
  } ;%
}


% Copyright (C) 2016  Joseph Rabinoff

% ipe2tikz is free software; you can redistribute it and/or modify it under
% the terms of the GNU General Public License as published by the Free
% Software Foundation; either version 3 of the License, or (at your option)
% any later version.

% ipe2tikz is distributed in the hope that it will be useful, but WITHOUT ANY
% WARRANTY; without even the implied warranty of MERCHANTABILITY or FITNESS
% FOR A PARTICULAR PURPOSE.  See the GNU General Public License for more
% details.

% You should have received a copy of the GNU General Public License along with
% ipe2tikz; if not, you can find it at "http://www.gnu.org/copyleft/gpl.html",
% or write to the Free Software Foundation, Inc., 675 Mass Ave, Cambridge, MA
% 02139, USA.


% ipe compatibility TikZ styles

\usetikzlibrary{arrows.meta}

\makeatletter

% These should behave almost exactly like ipe arrows.  They disable correcting
% for the miter length and line width.  This is important for visual consistency
% with ipe, since ipe arrows get much larger when the line width is increased.
% They also use the line join and cap styles from the main path.  These are very
% simple arrows: there is no harpoon version, and the convex hull computation is
% sloppy.

\pgfdeclarearrow{
  name = ipe _linear,
  defaults = {
    length = +1bp,
    width  = +.666bp,
    line width = +0pt 1,
  },
  setup code = {
    % Control points
    \pgfarrowssetbackend{0pt}
    \pgfarrowssetvisualbackend{
      \pgfarrowlength\advance\pgf@x by-.5\pgfarrowlinewidth}
    \pgfarrowssetlineend{\pgfarrowlength}
    \ifpgfarrowreversed
      \pgfarrowssetlineend{\pgfarrowlength\advance\pgf@x by-.5\pgfarrowlinewidth}
    \fi
    \pgfarrowssettipend{\pgfarrowlength}
    % Convex hull
    \pgfarrowshullpoint{\pgfarrowlength}{0pt}
    \pgfarrowsupperhullpoint{0pt}{.5\pgfarrowwidth}
    % The following are needed in the code:
    \pgfarrowssavethe\pgfarrowlinewidth
    \pgfarrowssavethe\pgfarrowlength
    \pgfarrowssavethe\pgfarrowwidth
  },
  drawing code = {
    \pgfsetdash{}{+0pt}
    \ifdim\pgfarrowlinewidth=\pgflinewidth\else\pgfsetlinewidth{+\pgfarrowlinewidth}\fi
    \pgfpathmoveto{\pgfqpoint{0pt}{.5\pgfarrowwidth}}
    \pgfpathlineto{\pgfqpoint{\pgfarrowlength}{0pt}}
    \pgfpathlineto{\pgfqpoint{0pt}{-.5\pgfarrowwidth}}
    \pgfusepathqstroke
  },
  parameters = {
    \the\pgfarrowlinewidth,%
    \the\pgfarrowlength,%
    \the\pgfarrowwidth,%
  },
}


\pgfdeclarearrow{
  name = ipe _pointed,
  defaults = {
    length = +1bp,
    width  = +.666bp,
    inset  = +.2bp,
    line width = +0pt 1,
  },
  setup code = {
    % Control points
    \pgfarrowssetbackend{0pt}
    \pgfarrowssetvisualbackend{\pgfarrowinset}
    \pgfarrowssetlineend{\pgfarrowinset}
    \ifpgfarrowreversed
      \pgfarrowssetlineend{\pgfarrowlength}
    \fi
    \pgfarrowssettipend{\pgfarrowlength}
    % Convex hull
    \pgfarrowshullpoint{\pgfarrowlength}{0pt}
    \pgfarrowsupperhullpoint{0pt}{.5\pgfarrowwidth}
    \pgfarrowshullpoint{\pgfarrowinset}{0pt}
    % The following are needed in the code:
    \pgfarrowssavethe\pgfarrowinset
    \pgfarrowssavethe\pgfarrowlinewidth
    \pgfarrowssavethe\pgfarrowlength
    \pgfarrowssavethe\pgfarrowwidth
  },
  drawing code = {
    \pgfsetdash{}{+0pt}
    \ifdim\pgfarrowlinewidth=\pgflinewidth\else\pgfsetlinewidth{+\pgfarrowlinewidth}\fi
    \pgfpathmoveto{\pgfqpoint{\pgfarrowlength}{0pt}}
    \pgfpathlineto{\pgfqpoint{0pt}{.5\pgfarrowwidth}}
    \pgfpathlineto{\pgfqpoint{\pgfarrowinset}{0pt}}
    \pgfpathlineto{\pgfqpoint{0pt}{-.5\pgfarrowwidth}}
    \pgfpathclose
    \ifpgfarrowopen
      \pgfusepathqstroke
    \else
      \ifdim\pgfarrowlinewidth>0pt\pgfusepathqfillstroke\else\pgfusepathqfill\fi
    \fi
  },
  parameters = {
    \the\pgfarrowlinewidth,%
    \the\pgfarrowlength,%
    \the\pgfarrowwidth,%
    \the\pgfarrowinset,%
    \ifpgfarrowopen o\fi%
  },
}


% For correcting minipage width in stretched nodes
\newdimen\ipeminipagewidth
\def\ipestretchwidth#1{%
  \pgfmathsetlength{\ipeminipagewidth}{#1/\ipenodestretch}}

\tikzstyle{ipe import} = [
  % General ipe defaults
  x=1bp, y=1bp,
%
  % Nodes
  ipe node stretch/.store in=\ipenodestretch,
  ipe stretch normal/.style={ipe node stretch=1},
  ipe stretch normal,
  ipe node/.style={
    anchor=base west, inner sep=0, outer sep=0, scale=\ipenodestretch
  },
%
  % Use a special key for the mark scale, so that the default can be overriden.
  % (This doesn't happen with the scale= key; those accumulate.)
  ipe mark scale/.store in=\ipemarkscale,
%
  ipe mark tiny/.style={ipe mark scale=1.1},
  ipe mark small/.style={ipe mark scale=2},
  ipe mark normal/.style={ipe mark scale=3},
  ipe mark large/.style={ipe mark scale=5},
%
  ipe mark normal, % Set default
%
  ipe circle/.pic={
    \draw[line width=0.2*\ipemarkscale]
      (0,0) circle[radius=0.5*\ipemarkscale];
    \coordinate () at (0,0);
  },
  ipe disk/.pic={
    \fill (0,0) circle[radius=0.6*\ipemarkscale];
    \coordinate () at (0,0);
  },
  ipe fdisk/.pic={
    \filldraw[line width=0.2*\ipemarkscale]
      (0,0) circle[radius=0.5*\ipemarkscale];
    \coordinate () at (0,0);
  },
  ipe box/.pic={
    \draw[line width=0.2*\ipemarkscale, line join=miter]
      (-.5*\ipemarkscale,-.5*\ipemarkscale) rectangle
      ( .5*\ipemarkscale, .5*\ipemarkscale);
    \coordinate () at (0,0);
  },
  ipe square/.pic={
    \fill
      (-.6*\ipemarkscale,-.6*\ipemarkscale) rectangle
      ( .6*\ipemarkscale, .6*\ipemarkscale);
    \coordinate () at (0,0);
  },
  ipe fsquare/.pic={
    \filldraw[line width=0.2*\ipemarkscale, line join=miter]
      (-.5*\ipemarkscale,-.5*\ipemarkscale) rectangle
      ( .5*\ipemarkscale, .5*\ipemarkscale);
    \coordinate () at (0,0);
  },
  ipe cross/.pic={
    \draw[line width=0.2*\ipemarkscale, line cap=butt]
      (-.5*\ipemarkscale,-.5*\ipemarkscale) --
      ( .5*\ipemarkscale, .5*\ipemarkscale)
      (-.5*\ipemarkscale, .5*\ipemarkscale) --
      ( .5*\ipemarkscale,-.5*\ipemarkscale);
    \coordinate () at (0,0);
  },
%
  % Arrow sizes (for TikZ arrows)
  /pgf/arrow keys/.cd,
  ipe arrow normal/.style={scale=1},
  ipe arrow tiny/.style={scale=.4},
  ipe arrow small/.style={scale=.7},
  ipe arrow large/.style={scale=1.4},
  ipe arrow normal,
  /tikz/.cd,
%
  % Approximations to ipe arrows
  % Put in a style to allow to reset default scale when "ipe arrow normal" is
  % changed.  I think this is the only way, since all the parameters to arrows
  % are expanded when the tip is declared.
  ipe arrows/.style={
    ipe normal/.tip={
      ipe _pointed[length=1bp, width=.666bp, inset=0bp,
                   quick, ipe arrow normal]},
    ipe pointed/.tip={
      ipe _pointed[length=1bp, width=.666bp, inset=0.2bp,
                   quick, ipe arrow normal]},
    ipe linear/.tip={
      ipe _linear[length = 1bp, width=.666bp,
                  ipe arrow normal, quick]},
    ipe fnormal/.tip={ipe normal[fill=white]},
    ipe fpointed/.tip={ipe pointed[fill=white]},
    ipe double/.tip={ipe normal[] ipe normal},
    ipe fdouble/.tip={ipe fnormal[] ipe fnormal},
    % These should maybe use [bend], but that often looks bad unless it's on an
    % actual arc.
    ipe arc/.tip={ipe normal},
    ipe farc/.tip={ipe fnormal},
    ipe ptarc/.tip={ipe pointed},
    ipe fptarc/.tip={ipe fpointed},
  },
  ipe arrows, % Set default sizes
]

% I'm not sure how to do this in a .style, since the #args get confused.
\tikzset{
  rgb color/.code args={#1=#2}{%
    \definecolor{tempcolor-#1}{rgb}{#2}%
    \tikzset{#1=tempcolor-#1}%
  },
}

\makeatother

\endinput

\usetikzlibrary{matrix,positioning,decorations.pathreplacing}
\usetikzlibrary{calc,quotes,angles}
\usetikzlibrary{arrows, arrows.meta, patterns}

\usetikzlibrary{decorations.pathreplacing}
\tikzset{
    position label/.style={
       above = 3pt,
       text height = 2ex,
       text depth = 1ex
    }
}

% \usetikzlibrary{decorations.markings}
\tikzset{
  font={\fontsize{14pt}{12}\selectfont}
}



\useoutertheme[subsection=false,shadow]{miniframes}
\useinnertheme{default}
\usefonttheme{serif}
%\usepackage{palatino}
\usepackage{mathpazo}
%\usepackage{utopia}
\usepackage{stmaryrd} % for varodot, bigodot 
\usepackage{mathabx} % for \coAsterisk
%\usepackage{mnsymbol}
%\setbeamertemplate{itemize item}{\scriptsize\raise1.7pt\hbox{\donotcoloroutermaths$\Asterisk$}}
%\setbeamertemplate{itemize item}{\scriptsize\raise1.7pt\hbox{\donotcoloroutermaths$\varodot$}}
%\setbeamertemplate{itemize subitem}{\scriptsize\raise1.25pt\hbox{\donotcoloroutermaths$\rhd$}}

\usepackage{xifthen}% provides \isempty tesst

\setbeamerfont{title like}{shape=\scshape}
\setbeamerfont{frametitle}{}



\setbeamercolor*{lower separation line head}{bg=blue} 
\setbeamercolor*{normal text}{fg=black,bg=white} 
\setbeamercolor*{alerted text}{fg=red} 
\setbeamercolor*{example text}{fg=black} 
%\setbeamercolor*{frametitle}{fg=aimsbrown} 
\setbeamercolor*{frametitle}{fg=imperialBlue} 
\setbeamercolor*{structure}{fg=black} 
 
\setbeamercolor*{palette tertiary}{fg=black,bg=black!10} 
\setbeamercolor*{palette quaternary}{fg=black,bg=black!10} 

%\renewcommand{\(}{\begin{columns}}
%\renewcommand{\)}{\end{columns}}
%\newcommand{\<}[1]{\begin{column}{#1}}
%\renewcommand{\>}{\end{column}}

% ======================================
% custom commands 
\newcommand{\cemph}[1]{\textcolor{\highlightcolor}{#1}}
\newcommand{\calert}[1]{\textcolor{red}{#1}}

\setbeamertemplate{navigation symbols}{}
%\renewcommand\frametitle[1]{{\textsc{\Large \textcolor{\highlightcolor}{#1}}}\vspace{0.6cm}\par}

\setbeamertemplate{frametitle}
{
{\textsc\bf \insertframetitle}\vspace{0.2cm}\par
}


%%%%%%%%%%%%%%%%%%%%%%%%%%%%%%%%%%%%%%%%%%%%%%%%%%
\setbeamertemplate{headline}{% 
	\setbeamercolor{head1}{bg=\headbarcolor}
	 \hbox{%
  \begin{beamercolorbox}[wd=.01\paperwidth,ht=2.25ex,dp=50ex,center]{head1}%
  \fontsize{5}{5}\selectfont  
  \end{beamercolorbox}%
  }
  \vspace{-50ex}
}
\setbeamertemplate{footline}{
\begin{tiny}
\setbeamercolor{foot1}{fg=black,bg=gray!10}
\setbeamercolor{foot2}{fg=gray,bg=gray!15}
\setbeamercolor{foot3}{fg=gray,bg=gray!10}
\setbeamercolor{foot4}{fg=black,bg=gray!20}
\setbeamercolor{foot5}{fg=gray,bg=gray!15}
\setbeamercolor{foot6}{fg=black,bg=gray!20}

% taken from theme infolines and adapted
  \leavevmode%
  \hbox{%
  \begin{beamercolorbox}[wd=.45\paperwidth,ht=2.25ex,dp=1ex,center]{foot1}%
  \fontsize{5}{5}\selectfont
  \flushleft \hspace*{2ex}{\footertitle}
  \end{beamercolorbox}%
  % \begin{beamercolorbox}[wd=.08\paperwidth,ht=2.25ex,dp=1ex,center]{foot2}
  % \end{beamercolorbox}%
  %   \begin{beamercolorbox}[wd=.05\paperwidth,ht=2.25ex,dp=1ex,center]{foot3}
  % \end{beamercolorbox}%
    \begin{beamercolorbox}[wd=.45\paperwidth,ht=2.25ex,dp=1ex,center]{foot4}%
  \fontsize{5}{5}\selectfont
  \authorname\hspace{5mm}@\location, \talkDate%\ (\authorweb) 
  \end{beamercolorbox}%
  % \begin{beamercolorbox}[wd=.05\paperwidth,ht=2.25ex,dp=1ex,center]{foot5}
  % \end{beamercolorbox}%
  \begin{beamercolorbox}[wd=.1\paperwidth,ht=2.25ex,dp=1ex,right]{foot6}%
	\insertframenumber{}  \hspace*{2ex} 
  \end{beamercolorbox}}%
  \vskip0pt%
\end{tiny}
\vskip0pt
}


\setbeamertemplate{blocks}[rounded][shadow=false]


\newenvironment<>{myblock}[1]{%
  \begin{actionenv}#2%
      \def\insertblocktitle{#1}%
      \par%
      \mode<presentation>{%
%       \setbeamercolor{block title}{fg=black,bg=aimslightbrown!50!white}
      \setbeamercolor{block title}{fg=black,bg=imperialBlue!45!white}
       \setbeamercolor{block body}{fg=black,bg=gray!20}
       \setbeamercolor{itemize item}{fg=blue!40!white}
       \setbeamertemplate{itemize item}[triangle]
     }%
      \usebeamertemplate{block begin}}
    {\par\usebeamertemplate{block end}\end{actionenv}}

\newenvironment<>{myblock2}[1]{%
  \begin{actionenv}#2%
      \def\insertblocktitle{#1}%
      \par%
      \mode<presentation>{%
       \setbeamercolor{block title}{fg=white,bg=blue!80!black}
       \setbeamercolor{block body}{fg=black,bg=gray!20}
       \setbeamercolor{itemize item}{fg=green!60!black}
       \setbeamertemplate{itemize item}[triangle]
     }%
      \usebeamertemplate{block begin}}
    {\par\usebeamertemplate{block end}\end{actionenv}}

\gdef\colchar#1#2{%
  \tikz[baseline]{%
%  \node[anchor=base,inner sep=2pt,outer sep=0pt,fill = #2!20]
%  {\large{#1}};
  \node[anchor=base,inner sep=1pt,outer sep=0pt,fill = #2!20]
  {{\fontsize{11}{13}\selectfont #1}};
    }%
}%
\gdef\drawfontframe#1#2{%
  \tikz[baseline]{%
  \node[anchor=base,inner sep=2pt,outer sep=0pt,fill = #2!20] {#1};
    }%
  }%


\makeatletter
\let\@@magyar@captionfix\relax
\makeatother

%%% Local Variables:
%%% mode: latex
%%% TeX-master: "2018-09-arusha-linear-regression"
%%% End:





\newif\iflattersubsect

\AtBeginSection[] {
    \begin{frame}<beamer>
    \frametitle{Overview} %
    \tableofcontents[currentsection]  
    \end{frame}
    \lattersubsectfalse
}

\AtBeginSubsection[] {
    \iflattersubsect
    \begin{frame}<Coming Next>
    \frametitle{Overview} %
    \tableofcontents[currentsubsection]  
    \end{frame}
    \fi
    \lattersubsecttrue
}

\begin{document}


%%%%%%%%%%%%%%%%%%%%%%%%%%%%%%%%%%%%%%%%%%%%%%%%%%%%%%

{\setbeamertemplate{footline}{}
\begin{frame}
\title{\slidesettitle}
%\subtitle{SUBTITLE}
\author{\footnotesize
  \textbf{\authorname}
 }

 %%% LOGO

% \begin{flushright}
%   % \begin{columns}
%   %   \column{0.5\hsize}
%   %   \column{0.45\hsize}
%\includegraphics[height = 8mm]{./figures/qla}\hspace{2mm}
%     \includegraphics[height = 8mm]{./figures/aims-rwanda}\\[2mm]
%\includegraphics[height = 8mm]{./figures/imperial}
%%\end{columns}
%\end{flushright}

\vspace{-0cm}
%\begin{flushleft}
%\vspace{-1.5cm}{\small \textcolor{blue}{\coursetitle}}\\\vspace{2cm}
{\huge \slidesettitle \ifthenelse{\equal{\slidesetsubtitle}{}}%
    {}% if #1 is empty
    {: \\ {\large \slidesetsubtitle}}% if #1 is not empty
    } \\    
    %\vspace{20pt}
%\end{flushleft}
  
 
% this is all stuff below the talk title. make two columns, just in
% case you want to have a picture or a second affiliation here 
\begin{columns}[t]
\column{0.8\hsize}
%\begin{flushleft}
\begin{columns}[t]
\column{0.6\hsize}
\insertauthor \\[2mm]
\authoraffiliation\\[2mm]
\column{0.25\hsize}
\\[2mm]
\includegraphics[height = 0.3cm]{./figures-general/twitter}{\small @\authortwitter}\\[-1mm]
\mbox{\small \url{\authoremail}}
\end{columns}
\column{0.14\hsize}
\end{columns}
% \authorweb\\
\vspace{7mm}
% \aimslightbrown{The Nelson Mandela African Institute of Science and
%   Technology\\Arusha, Tanzania}\\[2mm]
\insertdate
%\end{flushleft}
\end{frame}
}

%%% Local Variables:
%%% mode: latex
%%% TeX-master: t
%%% End:

\linespread{1.2} 

\makeatletter
\newcommand{\Pause}[1][]{\unless\ifmeasuring@\relax
\pause[#1]%
\fi}
\makeatother







\section{Index Notation}

\begin{frame}{Vector Differentiation: Index Notation}
Back to our linear regression problem:
\begin{equation}
L(\vtheta) = \sum_{n=1}^N (y_n - \vphi(x_n)\transpose\vtheta)^2 = ||\vy - \Phi(X)\vtheta||^2
\end{equation}

How to find $\pderiv[L]{\vtheta} = [\pderiv[L]{\theta_i}]_i$? \pause

\begin{itemize}
\item Vector expressions can always be written in terms of operations on scalars. \pause
\item We know how to find partial derivatives of scalar expressions $\implies$ rewrite into index notation. \pause
\item Take partial derivatives using tricks.
\end{itemize}
\end{frame}


\begin{frame}{Index Notation of Vector Expressions}
\begin{align}
\norm{\vx}^2 &= \sum_i x_i^2 \\ \Pause
\vx\transpose\vy &= \sum_i x_i y_i \\ \Pause
\vy = f(\vx) \vx &\implies y_i = f(\vx) x_i \\ \Pause
\mat A \vx &=\, ? \\ \Pause
{[\mat{A} \vx]_i} &= \sum_{j} A_{ij} x_j \\ \Pause
\mat A\vx &= \left[\sum_j A_{ij} x_j\right]_i
\end{align}
\end{frame}


\begin{frame}{Index Notation of Vector Expressions}
Back to our linear regression problem:
\begin{align}
L(\vtheta) &= \sum_{n=1}^N (y_n - \vphi(x_n)\transpose\vtheta)^2 = ||\vy - \Phi(X)\vtheta||^2 \\
&= \sum_{n=1}^N \left(y_n - \sum_{m=1}^M\phi_m(x_n) \theta_m\right)^2
\end{align}

We want to find $\pderiv[L]{\vtheta}$, i.e.~$[\pderiv[L]{\theta_i}]_i$ for all $i$.\pause
\begin{align}
&\pderiv{\theta_i} \sum_{n=1}^N \left(y_n - \sum_{m=1}^M\phi_m(x_n) \theta_m\right)^2 \\ \Pause
 &= \sum_n 2\left(y_n - \sum_{m=1}^M\phi_m(x_n)\theta_m\right) \pderiv{\theta_i} \left(y_n - \sum_{k=1}^M\phi_k(x_n) \theta_k\right)
\end{align}
\end{frame}

\begin{frame}{Index Differentiation}
\begin{align}
\pderiv[L]{\theta_i} &=  -2\sum_{n} \left(y_n - \sum_{m=1}^M\phi_m(x_n)\theta_m\right) \sum_k \phi_k(x_n) \pderiv[\theta_k]{\theta_i}
\end{align} \pause
\begin{itemize}
\item Think about $\pderiv[\theta_k]{\theta_i}$. What value does it take?\pause
\item If $i\neq k$, then $0$; if $i = k$, then $1$. \pause
\item Denoted as \emph{kronecker delta} $\delta_{ik}$. \pause
\item Think about what happens to $w_i = \sum_k \delta_{ik}$. \pause
\item Think about what happens to $x_k = \sum_a v_a \delta_{ak}$. \pause Is $\delta_{ak}$ a matrix? \pause
\end{itemize}
\begin{align}
\pderiv[L]{\theta_i}  &= -2\sum_{n} \left(y_n - \sum_{m=1}^M\phi_m(x_n)\theta_m\right) \sum_k \phi_k(x_n) \delta_{ki} \\ \Pause
&= -2\sum_{n} \left(y_n - \sum_{m=1}^M\phi_m(x_n)\theta_m\right)\phi_i(x_n)
\end{align}
\end{frame}


\begin{frame}{Vector Notation of Index Expressions}
\begin{align}
\pderiv[L]{\theta_i}  &= -2\sum_{n} \left(y_n - \sum_{m=1}^M\phi_m(x_n)\theta_m\right)\phi_i(x_n)
\end{align}

Sometimes we prefer a compact vector expression again. Index notation is ambiguous about row/col vector. If we remember convention, we can go back:
\begin{align}
&\pderiv[L]{\vtheta} = \uncover<2->{2\sum_n (\vphi(x_n)\transpose\vtheta-y_n)\vphi(x_n)\transpose && \text{Remember: } \vphi(x_n) \in \Reals^M} \nonumber \\
&= \onslide<3->{2\sum_n [\Phi(X)\vtheta - \vy]_n \vphi(x_n)\transpose && \text{Define: } \Phi(X)\!\in\!\Reals^{N\times M}, \Phi_{n:}(X) = \vphi(\vx_n)\transpose} \nonumber \\
&= \onslide<4->{2(\Phi(X)\vtheta - \vy)\transpose \Phi(X)} \nonumber\\
&= \onslide<5->{2(\vtheta\transpose\Phi(X)\transpose\Phi(X) - \vy\transpose\Phi(X))}
\end{align}
\end{frame}


\section{Differentiation of vector-valued functions}

\begin{frame}
  \frametitle{Chain Rule}
Back to our linear regression problem:
\begin{equation}
L(\vtheta) = \sum_{n=1}^N (y_n - \vphi(x_n)\transpose\vtheta)^2 = ||\vy - \Phi(X)\vtheta||^2
\end{equation}

\pause
\begin{itemize}
\item We could manually take partial derivatives of $L(\vtheta)$ \pause
\item Or, we could see our scalar function $L(\vtheta)$ as a function composition:
\end{itemize}
\begin{align*}
f(\vec g(\vtheta)) && f: \Reals^{D}\to\Reals \,, \qquad \vec g: \Reals^{E}\to\Reals^D
\end{align*}
\pause
\vspace{-0.4cm}
\begin{itemize}
\item Wouldn't it be nice to have a chain rule? \pause $\diffF{f(\vec g(\vtheta))}{\vtheta} = \diffF{f}{\vec g}\diffF{\vec g}{\vtheta}$ \pause
\item But how does a derivative of a vector w.r.t.~a vector work?
\end{itemize}
\end{frame}


\begin{frame}{Differentiation of a vector w.r.t.~a scalar}
\begin{itemize}
\item Differentiation quantifies how an output changes, \\
\hspace{4.3cm} in response to an input change. \pause
\item Change is measured by subtraction. {\color{gray} Draw vector subtraction.} \pause
\item Can use the same limiting argument as with scalars: \pause
\begin{align}
\deriv[\vx]{t} = \lim_{\Delta t\to 0} \frac{\vx(t + \Delta t) - \vx(t)}{\Delta t}
\end{align}
But now, derivative is a vector! \pause
\item Subtraction is elementwise, so we can reduce to scalar case:
\begin{align}
\left[\deriv[\vx]{t}\right]_i = \lim_{\Delta t\to 0} \frac{x_i(t + \Delta t) - x_i(t)}{\Delta t}
\end{align}
\end{itemize}
\end{frame}


\begin{frame}{Example: Circle}
Function that describes a point going round a circle, with period 1s:
\begin{align}
\vx = \begin{bmatrix}\cos 2\pi t & \sin 2\pi t\end{bmatrix}\transpose
\end{align}
{\color{gray} Draw on board} \pause

\begin{itemize}
\item Based on intuition, what should the speed be? \pause $2\pi$ m/s. \pause
\item Find velocity vector $\calcd\vx / \calcd t$.\pause
\item Find speed from the norm.
\end{itemize}
\end{frame}


\begin{frame}{Summary: Differentiation of a vector w.r.t.~a scalar}
\begin{myblock}{Derivative of a vector}
The derivative of a vector-valued function is given by the derivative of each of its elements.
\end{myblock}
\end{frame}







\section{Multivariate Chain Rule}
\begin{frame}[t]{Multivariate Chain Rule w.r.t.~scalar}
It turns out, there is a multivariate chain rule:
\only<1>{
\begin{align*}
\deriv[f(a(t), b(t))]{t} = \pderiv[f]{a}\deriv[a]{t} + \pderiv[f]{b}\deriv[b]{t}
\end{align*}}
\only<2->{
\begin{align*}
\deriv[f(\vec g(t))]{t} = \sum_{i=1}^D\pderiv[f]{g_i}\deriv[g_i]{t} && \vec g(t) \in \Reals^D
\end{align*}} \pause \pause
Inner product!  \pause \emph{Given our convention}, we can write in vector form:
\begin{align*}
\deriv[f(\vec g(t))]{t} = \underbrace{\deriv[f]{\vec g}}_{\text{row}}\cdot\underbrace{\deriv[\vec g]{t}}_{\text{column}}
\end{align*} \pause
\vspace{-0.3cm}
\begin{itemize}
\item $\deriv[\vec g]{t}$ is the derivative of a column vector. We keep this to be a column vector.
\item Can also be derived from a limit argument, like the scalar derivative (\emph{board}: Circle Example).
\end{itemize}
\end{frame}


\begin{frame}
  \frametitle{Example: Chain Rule}
  \begin{itemize}
    \item 
  Consider $f:\R^2\to \R, \quad \vec x:\R \to \R^2$
  \begin{align*}
    f(\vec x)&= f(x_1, x_2) = x_1^2+2x_2\,,\\
    \vec x(t) &=
                \begin{bmatrix}
                  x_1(t)\\ x_2(t)
                \end{bmatrix} =
                \begin{bmatrix}
      \sin(t)\\ \cos (t)
    \end{bmatrix}
  \end{align*}
  \vspace{-2mm}
  \pause
\item What are the dimensions of \blue{$\diffF{f}{\vec x}$} and \red{$\diffF{ \vec x}{t}$}?\\
  \only<2>{\qquad \qquad \qquad \emph{Work it out with your neighbors}}
  \pause \qquad \qquad \qquad \qquad \qquad \blue{$1\times 2$} and \red{$2\times 1$}
\item Compute the gradient $\diffF{f}{t}$ using the chain rule:
  \pause
  \begin{align*}
    \diffF{ f}{t} = \blue{\diffF{f}{\vec x}} \red{\diffF{\vec x}{t}}
     &= \blue{\begin{bmatrix}
		\frac{\partial f}{\partial x_1} & 
		\frac{\partial f}{\partial x_2}
		\end{bmatrix}}
		\red{\begin{bmatrix}
    	\frac{\partial x_1}{\partial t}\\[2mm]
        \frac{\partial x_2}{\partial t}
     	\end{bmatrix}}
 	= \blue{\begin{bmatrix}
		2\sin t & 2
		\end{bmatrix}}
        \red{\begin{bmatrix}
        \cos t\\
        -\sin t
        \end{bmatrix}}\\
     &= 2\sin t\cos t - 2\sin t = 2\sin t(\cos t -1)
  \end{align*}
\end{itemize}
\end{frame}




\begin{frame}[t]{Derivative of vector w.r.t.~vector}
We saw the chain rule if we were differentiating w.r.t.~a scalar:
\only<1>{
\begin{align*}
\deriv[f(a(t), b(t))]{t} = \pderiv[f]{a}\deriv[a]{t} + \pderiv[f]{b}\deriv[b]{t}
\end{align*}}
\only<2->{
\begin{align*}
\deriv[f(\vec g(t))]{t} = \sum_{i=1}^D\pderiv[f]{g_i}\deriv[g_i]{t} = \underbrace{\deriv[f]{\vec g}}_{\text{row}}\cdot\underbrace{\deriv[\vec g]{t}}_{\text{column}} && \vec g(t) \in \Reals^D
\end{align*}} \pause \pause
What happens if we differentiate w.r.t.~a vector? \pause \\
$\implies$ As before, we just stack the derivatives w.r.t.~each of the inputs.
\begin{align*}
\pderiv[f(\vec g(\vx))]{x_j} = \sum_{i=1}^D\pderiv[f]{g_i}\pderiv[g_i]{x_j} && \vec g(\vx) \in \Reals^D
\end{align*}
\end{frame}





\begin{frame}{Multivariate Chain Rule w.r.t.~vector}
This is a matrix multiplication! Can write in vector form:
\begin{align*}
\deriv[f(\vec g(\vx))]{\vx} = \underbrace{\deriv[f]{\vec g}}_{\text{row}}\cdot\underbrace{\deriv[\vec g]{\vx}}_{\text{matrix}}
\end{align*} \pause
\begin{itemize}
\item $\deriv[\vec g]{\vx}$ is the derivative of a column vector w.r.t.~the input vector $\vx$. We put the elements of $\vec g$ (i.e.~i) along the column, and the dimensions of the derivative (i.e.~j) along the row. \pause
\item Can also be derived from a directional derivative argument, but for a vector.
\end{itemize}
\end{frame}



\begin{frame}
  \frametitle{Vector Field Differentiation $f: \R^N\to\R^M$}
  \begin{align*}
    \vec y &= \vec f(\vec x) \in\R^M\,,\quad \vec x \in \R^N\\
    \begin{bmatrix}
     \colchar{$ y_1$}{red}\\
     \vdots\\
     \colchar{$y_M$}{blue}
    \end{bmatrix}
    &=
    \begin{bmatrix}
      \colchar{$f_1(\vec x)$}{red}\\
      \vdots\\
      \colchar{$f_M(\vec x)$}{blue}
    \end{bmatrix}
    =
    \begin{bmatrix}
      \colchar{$f_1(x_1,\dotsc, x_N)$}{red}\\
      \vdots\\
      \colchar{$f_M(x_1,\dotsc, x_N)$}{blue}
    \end{bmatrix}
  \end{align*}
  
  \pause
  \begin{itemize}    
    \item \cemph{Jacobian} matrix (collection of all partial derivatives)
      \begin{align*}
        \begin{bmatrix}
          \colchar{$\frac{d y_1}{d\vec x}$}{red}\\
          \vdots\\
          \colchar{$\frac{d y_M}{d\vec x}$}{blue}
        \end{bmatrix}
        =
        \begin{bmatrix}
          \colchar{$\frac{\partial f_1}{\partial
              x_1}$}{red} &
          \cdots & \colchar{$\frac{\partial
              f_1}{\partial x_N}$}{red}\\
          \vdots &  & \vdots\\
          \colchar{$\frac{\partial f_M}{\partial x_1}$}{blue} & \cdots & \colchar{$\frac{\partial
            f_M}{\partial x_N}$}{blue}
        \end{bmatrix}
                                                       \in\R^{M\times N}
      \end{align*}
    \end{itemize}
  
\end{frame}





\begin{frame}{Dimensionality of the Gradient}
  In general: A function $\vec
  f:\R^\text{\colchar{\scriptsize{$N$}}{blue}}\to\R^\text{\colchar{\scriptsize{$M$}}{orange}}$
  has a gradient that is an
  $M\times N$-matrix
  with
    $$
    \diffF{\vec f}{\vec x}\in \R^{M\times N}\,,\qquad \mathrm{d}\vec f[m,n] =
    \frac{\partial f_m}{\partial x_n}
    $$
    \begin{center}
     Gradient dimension: \colchar{ \# target dimensions }{orange} $\times$ \colchar{\# input dimensions}{blue}
    \end{center}

\pause
A function composition $\vf(\vx) = \vf_{\vg}(\vg(\vx))$ {\tiny(subscript to distinguish from overall function)} has the constraint that the \emph{output dimension} of $\mathbf g(\cdot)$ has to equal the \emph{input dimension} of $\vf_{\vg}(\cdot)$, so that we can compute $\vf(\mathbf g(\vx))$. \pause

This ensures that the shapes of the chain rule work out:
\begin{gather}
\vf : \Reals^\text{\colchar{\scriptsize{$N$}}{blue}} \to \Reals^\text{\colchar{\scriptsize{$M$}}{orange}} \qquad \qquad \vf_{\vg} : \Reals^\text{\colchar{\scriptsize{$L$}}{green}} \to \Reals^\text{\colchar{\scriptsize{$M$}}{orange}} \qquad \qquad \mathbf g : \Reals^\text{\colchar{\scriptsize{$N$}}{blue}} \to \Reals^\text{\colchar{\scriptsize{$L$}}{green}} \\
\underbrace{\diffF{\vf}{\vx}}_{\text{\colchar{\scriptsize{$M$}}{orange}}\times\text{\colchar{\scriptsize{$N$}}{blue}}} = \underbrace{\diffF{\vf_{\vg}}{\mathbf g}}_{\text{\colchar{\scriptsize{$M$}}{orange}}\times\text{\colchar{\scriptsize{$L$}}{green}}} \underbrace{\diffF{\mathbf g}{\vx}}_{\text{\colchar{\scriptsize{$L$}}{green}}\times\text{\colchar{\scriptsize{$N$}}{blue}}}
\end{gather}

    
\end{frame}




\begin{frame}
\frametitle{Example: Vector Field Differentiation}

\vspace{-5mm}
\begin{align*}
\vec f(\vec x) = \mat A\vec x\,,\qquad \vec f(\vec x)\in\R^M, \quad\mat A\in\R^{M\times
  N}, \quad \vec x\in\R^N
\end{align*}

\vspace{-5mm}
{\scriptsize 
		\begin{align*}
        \begin{bmatrix}
          \colchar{$y_1$}{red}\\
          \vdots\\
          \colchar{$y_M$}{blue}
        \end{bmatrix}
        =
        \begin{bmatrix}
          \colchar{$f_1(\vec x)$}{red}\\
          \vdots\\
          \colchar{$f_M(\vec x)$}{blue}
        \end{bmatrix}
        =
        \begin{bmatrix}
          \colchar{$A_{11}x_1$}{red} + \colchar{$A_{12}x_2$}{red} + &
          \cdots & + \colchar{$A_{1N}x_N$}{red}\\
          \vdots \qquad \qquad \vdots & \vdots & \quad  \vdots \\
          \colchar{$A_{M1}x_1$}{blue} + \colchar{$A_{M2}x_2$}{blue} + & 
          \cdots & 
          + \colchar{$A_{MN}x_N$}{blue}
        \end{bmatrix}
      \end{align*}}
\vspace{-1mm}
\begin{itemize}
\item Compute the gradient $\diffF{\vec f}{\vec x}$
  \begin{itemize}
  %\item Dimension of $\diffF{\vec f}{\vec x}$:
  %  \pause Since $\vec f:\R^N\to \R^M$ it follows that
  %  $\diffF{\vec f}{\vec x}\in\R^{M\times N}$
    \pause
  \item Gradient:
  	\vspace{-1mm}
    \begin{align*}
      f_{i}(\vec x) &= \sum_{k=1}^NA_{ik}x_k \qquad 
      \implies \frac{\partial f_i}{\partial x_j} =
              \sum_{k}A_{ik}\pderiv[x_k]{x_j} = \sum_{k}A_{ik}\delta_{kj} = A_{ij} \nonumber\\
      \onslide+<4>{
      \implies \diffF{\vec f}{\vec x} &= \begin{bmatrix}
        \frac{\partial f_1}{\partial x_1} & \cdots & \frac{\partial
          f_1}{\partial x_N}\\
        \vdots & & \vdots\\
        \frac{\partial f_M}{\partial x_1} & \cdots & \frac{\partial
          f_M}{\partial x_N}
      \end{bmatrix}
                                                     =
                                                     \begin{bmatrix}
A_{11} & \cdots & A_{1N}\\
\vdots & &\vdots\\
A_{M1} & \cdots & A_{MN}
\end{bmatrix}
                  =\mat A \in\R^{M\times N}}
    \end{align*}
  \end{itemize}
\end{itemize}
\end{frame}




%%%%%%%%%%%%%%%%%%%%%%%%%%%%%%%%%%%%%%%%%
\begin{frame}
  \frametitle{Example: Multivariate Chain Rule}
  \begin{itemize}
    \item Consider the function
  \begin{align*}
    L(\vec e) &= \tfrac{1}{2}\|\vec e\|^2 = \tfrac{1}{2}\vec e\T\vec e\\
    \vec e &= \vec y - \mat A\vec x\,,\quad \vec x\in\R^N\,,\mat
             A\in\R^{M\times N}\,,\vec e,\vec y\in\R^M
  \end{align*}
\item Compute the gradient $\diffF{L}{\vec x}$. What is the dimension/size of $\diffF{L}{\vec
  x}$?
  \only<1>{
  	\vspace{2mm}
	\begin{center}
	\textbf{Work it out with your neighbours}
	\end{center}}

  \pause
  \begin{align}
\diffF{L}{\vec x} &= \blue{\pdiffF{L}{\vec
                     e}}\red{\pdiffF{\vec e}{\vec
                     x}}\nonumber\\
    \blue{\pdiffF{L}{\vec e}} &\blue{=\vec e\T}\in
                                               \R^{1\times M}\,, \quad\quad \pderiv[L]{e_i} = \pderiv[]{e_i} \sum_{j}\frac{1}{2} e_j^2 = \sum_j \frac{1}{2} 2e_j \pderiv[e_j]{e_i} = e_i \nonumber \\
    \red{\pdiffF{\vec e}{\vec x}} &\red{=-\mat
                                                     A}\in\R^{M\times
                                                     N} \nonumber\\
    \implies \diffF{L}{\vec x} &=\blue{\vec e\T}(\red{-\mat A}) =
                                     \red{-}\blue{(\vec y - \mat
                                     A\vec x)\T} \red{\mat
                                     A}\in\R^{1\times N}\nonumber
  \end{align}
\end{itemize}
\end{frame}




%%%%%%%%%%%%%%%%%%%%%%%%%%%%%%%%%%%%%%%%%

\begin{frame}{Summary}
\begin{itemize}
\item Chain rule for multivariate functions
\item Derivatives of vectors w.r.t.~scalars.
\item Derivatives of vectors w.r.t.~vectors (and shapes).
\end{itemize}
\end{frame}





\end{document}
%%% Local Variables: 
%%% mode: latex
%%% TeX-master: t
%%% End: 
